\chapter{Detailed Algorithms and Pseudocode}

This appendix provides implementation-independent pseudocode for the main workflows. These algorithms can be used in Chapter 7 and Chapter 8 if the supervisor prefers a more formal presentation.

\section{Algorithm K.1 - MQTT sensor ingestion}

\begin{enumerate}

\item Receive MQTT message with topic and payload.

\item If topic is not sensors/data, route to device-status handler or ignore.

\item Parse payload as JSON.

\item Validate required fields such as device\_id and timestamp.

\item Normalize measurement names and units.

\item Resolve management zone from payload or device mapping.

\item Create missing sensor mappings if allowed.

\item Insert weather record when environmental fields exist.

\item Insert one sensor\_readings row per normalized measurement.

\item Update runtime MQTT status.

\item Trigger real-time Digital Twin update pipeline.

\end{enumerate}

\section{Algorithm K.2 - Real-time Digital Twin update}

\begin{enumerate}

\item Load latest readings and weather records for the zone.

\item Evaluate sensor freshness and health.

\item Compute current moisture, soil temperature, water balance, and stress status.

\item Build ML feature payload.

\item Call ML service /predict.

\item If ML succeeds, store prediction result and update AI fields.

\item If ML fails, compute fallback rule-based decision.

\item Upsert zone\_states.

\item Append zone\_state\_history.

\item Generate or update recommendation.

\item Expose updated state to dashboard.

\end{enumerate}

\section{Algorithm K.3 - What-if simulation}

\begin{enumerate}

\item Receive zone\_id, irrigation duration, projection horizon, and optional flow rate.

\item Load latest zone state and environment snapshot.

\item Create no-irrigation scenario by applying moisture decrease over horizon.

\item Create with-irrigation scenario by adding irrigation gain then applying evaporation/decay.

\item Select ML path according to recommended action.

\item Compute difference from live state, 24-hour projection, and 48-hour projection.

\item Store simulation scenario and result.

\item Return explanation and curves to frontend.

\end{enumerate}

\section{Algorithm K.4 - Recommendation validation}

\begin{enumerate}

\item User opens recommendation preview.

\item User validates recommendation.

\item Backend updates recommendation status to validated.

\item Backend creates control command and execution feedback placeholder if applicable.

\item Backend creates irrigation event for scheduled execution.

\item No direct MQTT command is sent unless user explicitly starts pump or event becomes due.

\end{enumerate}

\section{Algorithm K.5 - Scheduled irrigation execution}

\begin{enumerate}

\item Scheduler periodically queries irrigation\_events with due start time.

\item For each due event, publish ON\_D5 command to WeMos topic.

\item Mark event as running.

\item When end time is reached, publish OFF\_D5 command.

\item Mark event as executed.

\item Create execution\_feedback record.

\item Update dashboard state.

\end{enumerate}

\section{K.6 Formula notes for simple moisture projection}

The current simulation can be described qualitatively rather than as a final agronomic model. A simple projection combines current moisture, evaporation loss, irrigation gain, and safety bounds. In a future agronomic implementation, this should be replaced or calibrated using a soil water balance model including evapotranspiration, rainfall, infiltration, drainage, field capacity, wilting point, rooting depth, and crop coefficients.
